\begin{table}[H]
\centering
\caption{Nama Daerah Matriks Transisi Kategori IPM (UNDP) 2023 \(\rightarrow\) 2024} \label{tab:app1ipmname2023x2024}
\begin{adjustbox}{max width=\textwidth}
\begin{tabular}{|p{4cm}|p{10cm}|}
\hline
\textbf{Kategori} & \textbf{Daftar Wilayah} \\
\hline
Sangat Tinggi (2023) \(\rightarrow\) Sangat Tinggi (2024) & Kota Jakarta Selatan, Kota Jakarta Timur, Kota Jakarta Pusat, Kota Jakarta Barat, Kota Jakarta Utara, Kota Bandung, Kota Bekasi, Kota Depok, Kota Tangerang Selatan \\
\hline
Sangat Tinggi (2023) \(\rightarrow\) Tinggi (2024) & - \\
\hline
Sangat Tinggi (2023) \(\rightarrow\) Sedang (2024) & - \\
\hline
Sangat Tinggi (2023) \(\rightarrow\) Rendah (2024) & - \\
\hline
Tinggi (2023) \(\rightarrow\) Sangat Tinggi (2024) & Kota Cimahi, Kota Tangerang \\
\hline
Tinggi (2023) \(\rightarrow\) Tinggi (2024) & Kep. Seribu, Bogor, Bandung, Ciamis, Kuningan, Cirebon, Sumedang, Subang, Purwakarta, Karawang, Bekasi, Kota Bogor, Kota Sukabumi, Kota Cirebon, Kota Tasikmalaya, Kota Banjar, Tangerang, Kota Cilegon, Kota Serang \\
\hline
Tinggi (2023) \(\rightarrow\) Sedang (2024) & - \\
\hline
Tinggi (2023) \(\rightarrow\) Rendah (2024) & - \\
\hline
Sedang (2023) \(\rightarrow\) Sangat Tinggi (2024) & - \\
\hline
Sedang (2023) \(\rightarrow\) Tinggi (2024) & Bandung Barat \\
\hline
Sedang (2023) \(\rightarrow\) Sedang (2024) & Sukabumi, Cianjur, Garut, Tasikmalaya, Majalengka, Indramayu, Pangandaran, Pandeglang, Lebak, Serang \\
\hline
Sedang (2023) \(\rightarrow\) Rendah (2024) & - \\
\hline
Rendah (2023) \(\rightarrow\) Sangat Tinggi (2024) & - \\
\hline
Rendah (2023) \(\rightarrow\) Tinggi (2024) & - \\
\hline
Rendah (2023) \(\rightarrow\) Sedang (2024) & - \\
\hline
Rendah (2023) \(\rightarrow\) Rendah (2024) & - \\
\hline
\end{tabular}
\end{adjustbox}
\end{table}